\begin{topicbox}{Propositional Axiom Schemas}
\begin{exposition}
The propositional-logic development uses the following classical Hilbert-style
axiom schemas as logical starting points. They are schemas: each substitution
of formulas for the displayed schematic letters gives an axiom instance.
\end{exposition}
\end{topicbox}

\begin{axiombox}{Axiom Schema (Propositional Logic 1)}
\begin{axiom}[Propositional Axiom Schema 1]
\label{ax:propositional-logic-1}
For all formulas $P$ and $Q$,
\[
P \Longrightarrow (Q \Longrightarrow P).
\]
\end{axiom}
\end{axiombox}

\begin{remark*}[Interpretation]
If $P$ is available, then it remains available under any additional
assumption $Q$.
\end{remark*}

\NoLocalDependencies

\begin{axiombox}{Axiom Schema (Propositional Logic 2)}
\begin{axiom}[Propositional Axiom Schema 2]
\label{ax:propositional-logic-2}
For all formulas $P$, $Q$, and $R$,
\[
\bigl(P \Longrightarrow (Q \Longrightarrow R)\bigr)
\Longrightarrow
\bigl((P \Longrightarrow Q) \Longrightarrow (P \Longrightarrow R)\bigr).
\]
\end{axiom}
\end{axiombox}

\begin{remark*}[Interpretation]
Implications may be composed under a common hypothesis.
\end{remark*}

\NoLocalDependencies

\begin{axiombox}{Axiom Schema (Propositional Logic 3)}
\begin{axiom}[Propositional Axiom Schema 3]
\label{ax:propositional-logic-3}
For all formulas $P$ and $Q$,
\[
(\neg Q \Longrightarrow \neg P) \Longrightarrow (P \Longrightarrow Q).
\]
\end{axiom}
\end{axiombox}

\begin{remark*}[Interpretation]
This is the classical contrapositive principle. It is the axiom schema that
marks the background logic as classical rather than intuitionistic.
\end{remark*}

\NoLocalDependencies
